\chapter{Literature Review}
\label{chap:literature_review}
This chapter reviews the literature relevant to evaluating the robustness of automatic speech recognition (ASR) to overlapping speech. It is structured into three sections. The first section summarises established knowledge about multi-speaker datasets, synthetic mixture generation, overlap modelling, evaluation metrics and ASR architectures. The second section identifies the primary gaps motivating this dissertation. That includes the fact that existing resources do not jointly provide realistic overlap, factor-level control, and fair evaluation when comparing references and hypotheses with different structures. The third section explores recent work that has attempted to address these gaps through more realistic synthesis, more advanced pipelines and end-to-end ASR.

\section{Foundational Knowledge}
\subsection{Multi-Speaker Corpora for ASR Research}
Audio transcription research has depended on a relatively small set of corpora that preserve the acoustic and interactional realism of natural conversations. The \textit{AMI Meeting Corpus} is one of the foundational resources for multi-speaker ASR because it provides approximately 100 hours of monitored meetings with rich annotations and a variety of naturally occurring and scenario-driven interactions \cite{carletta2005ami}. The \textit{CHiME-6} challenge further extended that problem into distant multi-microphone, unsegmented recordings in home environments, introducing a benchmark where diarisation and recognition are performed under noise, reverberation and spontaneous conversational turn-taking \cite{watanabe2020chime6challengetacklingmultispeakerspeech}.

These datasets are foundational to ASR research because they retain the properties that are essential for downstream evaluation and application, like realistic microphone recording conditions, speaker interruptions, incomplete sentences and natural overlap. They together represent the trade-off between ecological validity and controlability. The real datasets preserve natural overlap patterns, acoustic effects and turn-taking, but they do not allow independent control of degradation factors. Inversely, synthetic datasets allow for precise control of factors but often lose the realism of natural conversation. This dissertation is designed to contribute to the middle ground, balancing realism and control.

Therefore, synthetic mixture generation has become a common technique in robust speech research. Earlier speech-separation datasets were often built by randomly mixing utterances under simplified conditions, but later datasets have progressively introduced realism by adding environmental noise, reverberation, corpus diversity and lexical continuity. \textit{WHAM!} added real ambient noise and argued that clean, fully overlapped benchmarks were too narrow for realistic evaluation \cite{wichern2019whamextendingspeechseparation}. \textit{WHAMR!} enhanced this by adding reverberation, modelling the spectral and temporal smearing that arises in indoor environments \cite{maciejewski2020whamrnoisyreverberantsinglechannel}. \textit{LibriCSS} then provided a more generalised alternative to wsj0-2mix-based datasets like \textit{WHAM!} with mixture derived from \textit{LibriSpeech} \cite{chen2020continuousspeechseparationdataset}. Finally, \textit{LibriMix} extended this by mixing in \text{VCTK} speech, increasing the total length and thereby making it more suitable for training purposes \cite{cosentino2020librimixopensourcedatasetgeneralizable}. The field has provided a rich set of resources for ASR research, but there is yet to be a framework that allows customised manipulation of overlap, noise, duration and speaker count in a reuseable way for generating synthetic mixtures.

\subsection{Overlap Modelling and Evaluation Metrics for Multi-Speaker ASR}
In this literature, the overlap ratio is a useful way to measure and control mixture difficulty. In LibriCSS, the overlap ratio is defined as the total duration of the overlapped region divided by the total speech duration \cite{chen2020continuousspeechseparationdataset}. This is methodologically important because natural conversation is not uniformly overlapped. It is often varied between different dialogue contexts \cite{cetin06_interspeech}. A benchmark that reduces these contexts into a single average condition hides what a model is actually robust to and what it is not. The evolution from a simple fully overlapped mixture to realistically partially overlapped speech signified an important progress in realism, showing the increasing emphasis on treating overlap as an experimental variable rather than a marginal condition \cite{yang2022simulatingrealisticspeechoverlaps}.

In ASR evaluation, the choice of metric can influence the weights of different aspects considered. In single-speaker ASR, word error rate (WER) is the standard metric, but in multi-speaker and overlapping settings, its calculation and interpretation are problematic. \citet{galibert2013methodologies} shows that classical implementations of WER and diarisation error rate (DER) that were designed for single-speaker speech need to be generalised and adapted to multi-speaker conditions to remain meaningful. Several recent works have formalised the metric for multi-speaker transcription in different aspects. \citet{von_Neumann_2023} shows that variants like cpWER or ORC-WER can be understood within a wider framework of long-form multi-speaker WER, while \citet{vonneumann2024meetevaltoolkitcomputationword} provides a toolkit for computing metrics like WER, cpWER, MIMO-WER and OCR-WER. \citet{von_Neumann_2023} also discusses the separation between speaker-attributed metrics and speaker-agnostic metrics by comparing asclite, OCR-WER, cpWER and MIMO-WER. The conclusion is that lexical recognition, temporal information and speaker attribution are all important factors in multi-speaker ASR evaluation, and that no single metric captures all of these aspects equally well. As a result, the choice of metrics must therefore be aligned with the application and the specific aspect of robustness being evaluated while methods of comparing mismatched reference and hypothesis structures remain underexplored.

\subsection{ASR Architectures and Benchmark Systems for Robust Speech Transcription}
The model architecture operating on top of these datasets and metrics has evolved substantially. One family of models is the self-supervised encoders like wav2vec-2.0, which show that pretraining on very large unlabeled examples can learn representations that transfer effectively to downstream ASR by fine-tuning on smaller labelled datasets \cite{baevski2020wav2vec20frameworkselfsupervised}. Another family is the large encoder-decoder models like Whisper, which uses weakly supervised pretraining and focuses on zero-shot transfer to improve robustness \cite{radford2022robustspeechrecognitionlargescale}. The third family is the conformer and transducer-based architectures like Parakeet that augmented the transformer with convolutional modules to capture both local patterns and global interactions and widely used in streaming ASR \cite{sekoyan2025canary1bv2parakeettdt06bv3efficient}. In modern multi-speaker ASR research, these models are often embedded inside complex pipelines that include front-end separation, speaker diarisation and transcription modules \citep{raj2020integrationspeechseparationdiarization, vonneumann2024meetingrecognitioncontinuousspeech}. In conclusion, the field already has mature corpora, mixture-generation, evaluation metrics and ASR models. However, the systematic framework for isolating the effects of overlap, noise and other factors under controlled environments remains less mature.

\section{Establishing a Gap in the Literature}
Despite the progress in multi-speaker ASR research, there are still important methodological gaps for this dissertation. Current real meeting corpora are ecologically valid, but they do not function as reusable testbeds in which overlap ratio, signal-to-noise ratio (SNR), clip duration, and speaker count can all be varied independently and systematically \citep{ami-corpus,watanabe2020chime6challengetacklingmultispeakerspeech}. On the other hand, synthetic datasets each solve only part of the problem. For example, WHAM! added noise, WHAMR! added reverberation, LibriMix provided more data and more realistic mixtures \citep{wichern2019whamextendingspeechseparation, maciejewski2020whamrnoisyreverberantsinglechannel, cosentino2020librimixopensourcedatasetgeneralizable}. However, none of these resources was designed primarily as a configurable framework for causal analysis of ASR robustness. Hence, they provided many excellent benchmarks, but less support for answering narrow experimental questions about the standalone and interactive effects of degradation factors on ASR performance.

A second gap concerns the evaluation under structural mismatch between references and hypotheses. Recent metric work has clarified how to judge speaker-attributed hypotheses, speaker-agnostic hypotheses and multi-output systems \citep{von_Neumann_2023, vonneumann2024meetevaltoolkitcomputationword,von_Neumann_2025}. However, in practice, many strong ASR systems still output a single sequential hypothesis even when the reference is multi-speaker, while overlap benchmark often encode reference speech as concurrent stream \citep{baevski2020wav2vec20frameworkselfsupervised,sekoyan2025canary1bv2parakeettdt06bv3efficient}. This creates an awkward in-between case. While cpWER and tcpWER are designed for multi-stream references, they assume a speaker-structured hypothesis. ORC-WER is designed to ignore speaker labels and isolate lexical correctness. However, what is not provided is an estimation algorithm tailored to the specific case of comparing a double-stream reference to a single-stream hypothesis.

This dissertation will therefore address two gaps. The first is the absence of a reusable, reconfigurable synthetic mixture generation framework for overlapping speech that can manipulate overlap, noise, duration and speaker count as experimental variables. The second is the absence of an estimation algorithm for the particular evaluation scenario in which a potentially overlapping double-stream reference must be compared with a single-stream hypothesis transcript. The proposed solution is motivated by the second gap. It is to make the evaluation of mismatched reference and hypothesis structure more informative and fair. In conclusion, these problems define the niche of this dissertation: controlled robustness benchmarking for overlapping speech and an evaluation method for mismatched reference and hypothesis structures.

\section{Recent Advances}
\subsection{Realism and Generalisation in Overlap Modelling}
Recent work has attempted to reduce some of these limitations. Some research has focused on improving the realism of synthetic overlap. \citet{yang2022simulatingrealisticspeechoverlaps} criticise the standard practice of randomly mixing utterances and propose a more realistic simulation by learning overlap patterns from real conversation, and representing overlap patterns as a sequence of discrete tokens. Their results show that the ASR model trained with a realistic overlap simulation consistently outperforms models trained with random overlap. In the benchmarking direction, \citet{cornell2024chime8dasrchallengegeneralizable} spurs the generalisation across speaker counts, formal vs informal settings, meeting duration, acoustic scenarios and recording configurations. Although they have made valuable progress from optimisation under specific conditions to broader generalisation, improving realism, they have yet to provide systematic frameworks to control overlap, SNR, duration, and speaker count in a clean design.

\subsection{Integrated Pipelines and End-to-End ASR for Overlapping Speech}
A second direction of work has concentrated on integrated meeting-recognition pipelines. \cite{raj2020integrationspeechseparationdiarization} shows that an end-to-end modular system that integrates independently trained speech separation, diarisation and transcription modules can achieve competitive performance on LibriCSS. \citet{chen2020continuousspeechseparationconformer} demonstrate that the combination of continuous speech separation (CSS) and conformer-based ASR can achieve state-of-the-art performance in both single-channel and multi-channel settings on LibriCSS. \citet{vonneumann2024meetevaltoolkitcomputationword} use a pipeline that combines the CSS system with TF-GridNet separation and a speaker-agnostic ASR and achieves outstanding performance on ORC-WER. They also propose a syntactically informed diarisation for turn detection, which results in exceptional performance on cpWER. The strength of these modules is that they treat meeting transcription as an all-in-one problem rather than a pure ASR problem. This, conversely, brings the drawback of ambiguity on the source of improvements, which makes it harder to locate the specific factors that contribute to the improvements and also makes it harder to locate the point of failure.

A third line of work has turned to an end-to-end ASR framework for speaker-attributed modelling. \citet{kanda2021endtoendspeakerattributedasrtransformer} present a transformer-based end-to-end model that performs speaker counting, speech recognition and speaker identification, reducing the dependence on separated separation and diarisation modules. \citet{chang2021hypothesisstitcherendtoendspeakerattributed} shows that such a model is susceptible to the mismatch between training and testing conditions and introduces a hypothesis stitcher that fuses multiple hypotheses into a single output to help improve accuracy. \citet{kanda2022transcribetodiarizeneuralspeakerdiarization} push this further by showing speaker-attributed transcription can be used to derive diarisation output, achieving a comparable diarisation error rate to existing diarisation methods and better if the number of speakers is not known. Conclusively, these studies are important because they treat lexical content and speaker identity as mutually informative independent tasks. On the flip side, they also reveal that since these systems are more specialised, the training assumptions are stronger, and the resulting outputs are less comparable to single-stream ASR baselines.

More recent papers have tried to reduce the gap between theoretical models and practical applications. \citet{cui2024improvingspeakerassignmentspeakerattributed} propose a pipeline for real-life applications that involves voice activity detection (VAD), speaker diarisation (SD), and speaker-attributed ASR (SA-ASR). They show that using VAD results to finetune the SA-ASR model significantly reduces the speaker error rate (SER) and demonstrate that extracting speaker embedding from SD can further improve the SER. \citet{raj2024speakerattributionsurt} push the SURT framework from speaker-agnostic to speaker-attributed ASR by adding a speaker branch and prefixing each chunk with speaker information. \citet{zheng2025dncasrendtoendtrainingspeakerattributed} introduces the DNCASR model for joint neural speaker clustering and ASR, effectively addresses overlapping speech and achieves strong cpWER performance on AMI. These studies indicate that the field is moving towards more application-oriented systems of long-form, speaker-consistent and meeting-level transcription. Nonetheless, they still rely on a small set of benchmarks. They therefore match the need for studies that evaluate ASR systems under different controlled conditions.

\subsection{Unified Metrics for Multi-Speaker ASR Evaluation}
Recent literature on ASR metrics provides evolving insights into multi-speaker evaluation. \citet{vonneumann2024meetevaltoolkitcomputationword} contribute to standardising the computation of WER, cpWER, ORC-WER and MIMO-WER as a toolkit. \citet{von_Neumann_2025} contributes to the conceptual framework for clarifying which metrics are suitable for which output structure and which aspect of performance. They also propose efficient approximation algorithms for expensive metrics like ORC-WER and DI-cpWER. These works discourage the default use of a single metric for outputs that embody different assumptions about speaker structure and timing. However, even this more mature scoring literature still treats the metrics primarily as ways of assessing matching systems. It does not directly solve the narrower experimental problem of how to compare double-stream references with single-stream hypotheses. That is, therefore, the space in which the present study makes its methodological contribution.

\subsection{Summary}
Overall, parallel research has advanced in three main directions. It has made the synthetic overlap more realistic rather than purely random. It has strengthened meeting transcription through integrated modular pipelines that combine separation, diarisation, and recognition. It has progressed ASR evaluation metrics towards a more comprehensive approach. However, it still leaves a gap in not having a lightweight, reusable framework for systematically varying degradation factors and the metrics that handle stream-mismatched evaluation. The present dissertation is designed to contribute to that missing middle ground.